% TOP_PROBLEM: {"id": 20001138, "title": "Exact subgroup self-sumsets and a robust quotient-character leakage bound", "category_id": 2, "category": {"id": 2, "name": "combinatorics", "display_name": "Combinatorics", "description": "Counting problems, graph theory, discrete structures.", "slug": "combinatorics", "order_index": 2, "created_at": "2026-07-31T15:26:25.671Z"}, "status": "partially_solved", "problem_number": "AIM-COMBINATORICS-0263", "set_id": 14, "statusQualification": "Uses the precise containment subquestion. The exact-equality part is resolved: Theorem 1.1 of Yip--Yoo v2 forces two-element summands and a four-element subgroup, impossible when the summands coincide. Does not impose an undefined notion of slightly larger."}
% FIRST_POSTED: 2026-10-01
% ENTRY_KIND: statement_only
% UnsolvedMath ID 20001138; source catalogue code AIM-COMBINATORICS-0263
% !TeX program = lualatex
% Definition and sources only; this file is not an LLM solution attempt.
\documentclass[11pt,a4paper]{article}
\usepackage[margin=25mm]{geometry}
\usepackage{fontspec}
\setmainfont{lmroman10-regular.otf}[BoldFont=lmroman10-bold.otf,
 ItalicFont=lmroman10-italic.otf,BoldItalicFont=lmroman10-bolditalic.otf]
\usepackage{amsmath,amssymb,mathtools}
\usepackage{unicode-math}
\setmathfont{latinmodern-math.otf}
\AtBeginDocument{\let\setminus\smallsetminus}
\usepackage{xurl,hyperref}
\hypersetup{colorlinks=true,urlcolor=blue}
\setcounter{secnumdepth}{0}
\setlength{\parindent}{0pt}
\setlength{\parskip}{5pt}
\setlength{\emergencystretch}{3em}
\begin{document}
\section*{Exact subgroup self-sumsets and a robust quotient-character leakage bound}\label{20001138}
{\small UnsolvedMath ID 20001138 · AIM-COMBINATORICS-0263 · Combinatorics · AIM Workshop Problem Lists}
\subsection{Definitions and mathematical statement}
For an odd prime $p$, let $\mathbb F_p^*$ be the multiplicative group of its nonzero elements.
A proper multiplicative subgroup $H<\mathbb F_p^*$ is required to differ from the whole multiplicative group.
For $A\subseteq\mathbb F_p$, write
$A+A=\{a+b:a,b\in A\}$, allowing equal summands.
The containment $A+A\subseteq H$ includes every such sum, so in particular it excludes zero from $A+A$.

Define
\[
 Q(p)=\max\{|A|:A\subseteq\mathbb F_p,
       \exists H<\mathbb F_p^*\text{ with }A+A\subseteq H\}.
\]
The maximum includes the empty set and is therefore always defined.
The explicit containment version of the workshop question asks whether
\[
 Q(p)=p^{o(1)}\qquad(p\to\infty\text{ through odd primes}),
\]
meaning that for every $\varepsilon>0$ there is $p_0(\varepsilon)$ such that $Q(p)\leq p^\varepsilon$ for all primes $p\geq p_0$.
Equivalently, can a fixed positive power of $p$ ever be exceeded along infinitely many primes by a set whose every pairwise sum lies in one proper multiplicative subgroup?
A negative answer to the proposed subpower bound requires such an infinite family for some fixed exponent.
The subgroup may depend on the prime and on $A$; uniformity over all proper subgroups is part of the assertion.
The source also asks about exact equality $A+A=H$. The classification in [S3] rules out that equality for $|A|>1$, while the containment target formulated here is a different remaining question.
\subsection{Short English statement}
Can a set of polynomial size in a prime field have every pairwise sum inside one proper multiplicative subgroup?
\subsection{Sources}
\begin{itemize}
\item[S1] ULAM.AI and the UnsolvedMath Contributors, supplied problems.json, record 20001138 (AIM-COMBINATORICS-0263), AIM Workshop Problem Lists. Catalogue identity and source transcription; CC BY 4.0.
\url{https://huggingface.co/datasets/ulamai/UnsolvedMath}.
\item[S2] American Institute of Mathematics, Recent trends in additive combinatorics, item 4.7. Original workshop question underlying AIM-COMBINATORICS-0263.
\url{https://aimath.org/WWN/additivecomb/additivecomb.pdf}.
\item[S3] Chi Hoi Yip and Semin Yoo, Additive decompositions of multiplicative subgroups via differential identities, Theorem 1.1. Resolution of the separate exact-equality question.
\url{https://arxiv.org/pdf/2608.02568v2}.
\end{itemize}
\end{document}
